% Part: first-order-logic
% Chapter: introduction
% Section: first-order-logic

\documentclass[../../../include/open-logic-section]{subfiles}

\begin{document}

\olfileid{fol}{int}{fol}

\olsection{మొదటిస్థాయి విధేయ తర్కం (First-Order Logic)}

ఔపచారిక తర్కానికి మీ మొదటి పరిచయం నుంచే మొదటిస్థాయి విధేయ తర్కం మీకు
బహుశా పరిచయమై ఉంటుంది.\footnote{నిజానికి, అది మీకు పరిచయమే అని
దాదాపుగా ఊహిస్తున్నాం!{} పరిచయం లేకపోతే, \emph{forall x}
\citep{Magnus2021} వంటి మరింత ప్రాథమిక పాఠ్యగ్రంథాన్ని పునశ్చరణ
చేయవచ్చు.} అది ``పరిమాణీకరణ తర్కం'' లేదా ``విధేయ తర్కం'' అనే పేర్లతో
మీకు తెలిసి ఉండవచ్చు. అన్నిటికన్నా ముందు, మొదటిస్థాయి విధేయ తర్కం ఒక
ఔపచారిక భాష. అంటే దానికి ఒక నిర్దిష్ట సంకేతావళి ఉంటుంది; దాని
వ్యక్తీకరణలు ఆ సంకేతావళితో ఏర్పడిన సంకేతమాలలు. అయితే ప్రతి
సంకేతమాలకూ అనుమతి లేదు. అనుమతించిన వ్యక్తీకరణలు పదాలు,
\tetoken{సూత్రాలు}{!!{formula}s}, \tetoken{వాక్యాలు}{!!{sentence}s} అనే వేర్వేరు రకాల్లో ఉంటాయి. మనకు ప్రధాన
ఆసక్తి మొదటిస్థాయి విధేయ తర్కపు \tetoken{వాక్యాలలోనే}{!!{sentence}s}: అవి సహజభాషా వాక్యాలకు
ఔపచారిక సాదృశ్యాలను ఇస్తాయి; వాటి గురించి తర్కవేత్తకు సాధారణంగా
ఆసక్తి కలిగించే ప్రశ్నలను అడగవచ్చు. ఉదాహరణకు:
\begin{itemize}
    \item $!A$ నుంచి $!B$ తార్కికంగా అనుగమిస్తుందా?
    \item $!A$ తార్కికంగా సత్యమా, తార్కికంగా అసత్యమా, లేక
    పరిస్థిత్యాధీనమా?
    \item $!A$, $!B$లు తుల్యాలా?
\end{itemize}

ఇవి ప్రధానంగా మొదటిస్థాయి విధేయ తర్కపు \tetoken{వాక్యాలు}{!!{sentence}s} యొక్క ``అర్థం'' గురించిన
ప్రశ్నలు. ఉదాహరణకు, $!A$ నుంచి $!B$ తార్కికంగా అనుగమిస్తుందా అనే
ప్రశ్నను ఒక తత్వవేత్త ఇలా విశ్లేషించవచ్చు: $!A$ సత్యమై $!B$ అసత్యమయ్యే
సందర్భం ఏదైనా ఉందా (అయితే $!B$, $!A$ నుంచి అనుగమించదు), లేక $!A$ను
సత్యం చేసే ప్రతి సందర్భమూ $!B$ను కూడా సత్యం చేస్తుందా (అయితే $!B$,
$!A$ నుంచి అనుగమిస్తుంది)? అయితే ``సందర్భం'' అంటే ఏమిటో మనకు ఇంకా
చెప్పలేదు---అది \emph{అర్థవిచారం} చేసే పని. మొదటిస్థాయి విధేయ తర్కపు
అర్థవిచారం, ``సందర్భం'' అనే తత్వవేత్త అంతర్బుద్ధికి గణితపరంగా
ఖచ్చితమైన నమూనాను ఇస్తుంది; అంతేకాక---ఇది ముఖ్యం---ఒక
\tetoken{ఒక వాక్యం}{!!a{sentence}}~$!A$ ఆ సందర్భంలో \emph{సత్యం కావడం} అంటే ఏమిటో కూడా
ఖచ్చితపరుస్తుంది. మనం అభివృద్ధి చేయబోయే ఆ గణితపరమైన నమూనాను
\tetoken{ఒక నిర్మాణం}{!!a{structure}} అంటాం. ``దానిలో సత్యం'' అనే భావనను ఖచ్చితపరిచే
సంబంధాన్ని \emph{సంతృప్తి} సంబంధం అంటాం. కాబట్టి \tetoken{వాక్యాలు}{!!{sentence}s}~$!A$,
\tetoken{నిర్మాణాలు}{!!{structure}s}~$\Struct{M}$ కోసం ``$\Struct{M}$లో $!A$ సంతృప్తి
చెందుతుంది'' (సంకేతాల్లో: $\Sat{M}{!A}$) అనే దాన్ని నిర్వచిస్తాం.
అది పూర్తయిన తరువాత ``నుంచి అనుగమిస్తుంది,'' ``తార్కికంగా సత్యం''
వంటి ఇతర అర్థపర పదాలకు కూడా ఖచ్చితమైన నిర్వచనాలు ఇవ్వవచ్చు.
ఉదాహరణకు,
$\lforall[x][(!A(x) \lif !B(x))], \lexists[x][!A(x)] \Entails
\lexists[x][!B(x)]$ అవుతుందో లేదో గణితపరమైన ఖచ్చితత్వంతో నిర్ణయించడానికి
ఈ నిర్వచనాలు వీలు కల్పిస్తాయి. సమాధానం సహజంగానే ``అవును.'' ఆంగ్ల
వాక్యాలను మొదటిస్థాయి విధేయ తర్కంలో సంకేతీకరించడంలో మీకు ఇప్పటికే శిక్షణ
ఉంటే, దీన్ని ఉదాహరణకు ``చీమలన్నీ కీటకాలు; చీమలు ఉన్నాయి; కాబట్టి
కీటకాలు ఉన్నాయి'' అనే వాదన సంకేతీకరణగా గుర్తిస్తారు. అది స్పష్టంగా
చెల్లుబాటైన వాదన; కాబట్టి మన ఔపచారిక భాషలోని ``నుంచి అనుగమిస్తుంది''
అనే భావనకు మన గణిత నమూనా కూడా అదే సమాధానం ఇవ్వాలి.
\sourcecorrection{OLTEINT-001}{ఆంగ్ల మూలంలో మొదటి పూర్వాధారంలోని
సార్వత్రిక పరిమాణీకరణ కుండలీకరణ, మొత్తం అనుగమన ప్రకటన చివరికి
జారిపోయింది. ఉదాహరణను వివరించే వెంటనే వచ్చే సహజభాషా వాదన
నిర్ణయించినట్లుగా, సార్వత్రిక పూర్వాధారం తరువాత కుండలీకరణను మూసి,
అనుగమన ప్రకటన చివరనున్న అదనపు కుండలీకరణను తొలగించాం.}

ఔపచారిక తర్కానికి మీ మొదటి పరిచయం నుంచి బహుశా గుర్తుండే మరో విషయం
\emph{\tetoken{వ్యుత్పత్తులు}{!!{derivation}s}} ఉండటం. మీరు మొదటి ఔపచారిక తర్క పాఠ్యాంశం
చదివి ఉంటే, మీ అధ్యాపకులు అలాంటి \tetoken{వ్యుత్పత్తులను}{!!{derivation}s} కనుగొనే సాధన
చేయించి ఉంటారు; పైనున్న అనుగమనం నిలుస్తుందని చూపే
\tetoken{ఒక వ్యుత్పత్తిని}{!!a{derivation}} కూడా చేయించి ఉండవచ్చు. \tetoken{వ్యుత్పత్తులను}{!!{derivation}s} ఇవ్వడానికి
అనేక పద్ధతులు ఉన్నాయి: మీరు ``సహజ నిగమనం'' లేదా ``సత్య వృక్షాలు''
అనే పద్ధతుల్లో ఏదో ఒకటి చేసి ఉండవచ్చు; మరెన్నో పద్ధతులు కూడా
ఉన్నాయి. పైనున్న తర్కవేత్తల ప్రశ్నలకు జవాబులు కనుగొనే సాధనాలను
అందించడమే \tetoken{వ్యుత్పత్తి}{!!{derivation}} వ్యవస్థల ఉద్దేశం. ఉదాహరణకు,
$\lforall[x][(!A(x) \lif !B(x))]$, $\lexists[x][!A(x)]$లు
పూర్వాధారాలుగా, $\lexists[x][!B(x)]$ నిష్కర్షగా (చివరి వరుసగా) ఉండే
సహజ నిగమన \tetoken{వ్యుత్పత్తి}{!!{derivation}}, $\lforall[x][(!A(x) \lif !B(x))]$,
$\lexists[x][!A(x)]$ల నుంచి $\lexists[x][!B(x)]$ తార్కికంగా
అనుగమిస్తుందని \emph{నిర్ధారిస్తుంది}.

అయితే అలా ఎందుకు? పైకి చూస్తే, \tetoken{వ్యుత్పత్తి}{!!{derivation}} వ్యవస్థలకు
అర్థవిచారంతో సంబంధమే లేదు: ఒక ఔపచారిక \tetoken{వ్యుత్పత్తి}{!!{derivation}} ఇవ్వడమంటే,
కొన్ని నియమాలకు లోబడి సంకేతాలను అమర్చడమే; అవి ``సందర్భాలు'' లేదా
``దానిలో సత్యం'' అనే విషయాలను అసలు ప్రస్తావించవు. \tetoken{వ్యుత్పత్తి}{!!{derivation}}
వ్యవస్థలకు, అర్థవిచారానికి మధ్య సంబంధాన్ని అధితార్కిక పరిశీలన ద్వారా
స్థాపించాలి. ఉదాహరణకు, $\lforall[x][(!A(x) \lif !B(x))]$,
$\lexists[x][!A(x)]$ అనే పూర్వాధారాల నుంచి $\lexists[x][!B(x)]$కు
  ఔపచారిక \tetoken{వ్యుత్పత్తి}{!!{derivation}} సాధ్యమవడం, ఆ రెండు పూర్వాధారాలు కలిపి
  $\lforall[x][(!A(x) \lif !B(x))]$, $\lexists[x][!A(x)]$ కలిపి
  $\lexists[x][!B(x)]$ను అనుగమింపజేసినప్పుడూ, అప్పుడు మాత్రమే అని
గణితపరంగా నిరూపించాలి. అయితే ఇది చేయడానికి ముందు, వాక్యనిర్మాణం,
అర్థవిచారం నిర్వచనాలను సరిగ్గా నిర్ధారించే చాలా శ్రమతో కూడిన పని
చేయాలి.
\sourcecorrection{OLTEINT-002}{ఆంగ్ల మూలంలోని రెండు నిగమన-ఉదాహరణ
ప్రస్తావనల్లో సార్వత్రిక పూర్వాధారాన్ని మూసే కుండలీకరణ కనిపించలేదు;
అదే సమయంలో నిష్కర్ష చివర అదనపు కుండలీకరణ కనిపించింది. ప్రతి
ప్రస్తావనలోనూ పూర్వాధారం, నిష్కర్షల ఉద్దేశిత పరిధులు స్పష్టంగా ఉండేలా
ఆ కుండలీకరణను సరైన స్థానానికి మార్చాం.}

\end{document}
